\section{2\quad 流体静力学基础}

%\subsection{流体平衡应力特征与静压强}

平衡下的流体应力特征：没有黏性切向力；只有法向力；不能承受拉力；压力是唯一表面力。
\pt{静止状态下}流体单位表面熵所处的压力称为压强，特征为：流体静压强方向与作用面\pt{垂直}，指向作用面\pt{内法线}方向；同一点熵的各方向流体静压强相等\m{p_n = p_x = p_y = p_z}。不同点的静压强是空间坐标的连续函数

%\subsection{流体平衡微分方程}

流体平衡微分方程：\m{\vec{f}-\frac{1}{\rho}\nabla p = 0}，单位质量质量力等于静压强合力。适用于静止或者相对静止状态的可压缩和不可压缩流体。\\
压强差公式：\m{\mathrm{d}p = \rho(f_x \mathrm{d}x+f_y \mathrm{d}y +f_z \mathrm{d}z)}。\\
流体平衡方程：\m{\left\{
    \begin{aligned}
         & \frac{\partial f_x}{\partial y}=\frac{\partial f_y}{\partial x} \\&\frac{\partial f_y}{\partial z}=\frac{\partial f_z}{\partial y}\\&\frac{\partial f_z}{\partial x}=\frac{\partial f_x}{\partial z}
    \end{aligned}
    \right.}\\
\pt{等压面}: 流体中压强相等的各点组成的面为等压面,等压面的势能也相等;等压面与质量力互相垂直;平衡状态，互不掺混的两种液体的分界面也是等压面。对于任意形式的连通器,在紧密相连而又属于同一性质的静止均值液体中深度相同的点压强相等。

%\subsection{流体静力学方程}

\pt{流体静力学基本方程}:\m{z+\frac{p}{\rho g}=c}.在重力作用下静止流体中各点的单位重量流体的总势能是相等的。其中 \m{z}为单位重量流体对某一基准面的\pt{位势能}，\m{p/\rho g}为单位重量流体的\pt{压强势能}，\m{c}为位势能和压强势能的总和,称为单位重量流体的\pt{总势能}。或表示为在重力作用下静止流体中各点的静水头都是相等的,则 \m{z}为单位重量流体的\pt{位置水头},\m{p/\rho g}为单位重量流体的\pt{压强水头},\m{c}为单位重量流体的\pt{静水头}。

%\subsection{压强的测量}

\pt{相对压强}:以当地大气压强为基准计算的压强称,\m{p_e = p - p-a}.
\pt{真空度}:容器内绝对压强小于当地大气压强时,\m{p_v = p_a - p}。

%\subsection{静止流体对物体表面的压强合力}

\pt{平面净水总压力}:\m{P = \rho g h_c A}压强合力等于受压面形心处的压强与面积的乘积,垂直平面壁\\
总压力的作用点 \m{z_d = z_c+\frac{J_{cx}}{z_c A}}

\pt{曲面净水总压力}:\pt{垂直分力}\m{P_z = \rho g V_P}和\pt{水平分力}\m{P_x = \rho g h_c A_x},总压力为\m{P = \sqrt{P_x^2 + P_z^2}}，总压力的作用点为\m{z_d = \frac{\int z \mathrm{d}P_z}{P_z}}。\pt{作用点}为 \m{P_x} 和 \m{P_z} 的作用线交点作 \m{\vec{P}} 的作用线与曲面交点。

%\subsection{例题分析}